\chapter{Cadence Matrix}
\label{chap:Cadence Matrix}

The Cadence Matrix below summarises the recommended review and inspection frequency for the principal assessment topic of each of the thirty chapters in this book. The matrix is intended as a practical planning tool for principal contractors, project managers, and health and safety advisers when preparing the health and safety monitoring programme for a construction project.

\textbf{How to use the matrix:} Identify the assessment topics relevant to your project by reviewing the chapter list. For each relevant topic, note the frequencies marked. ``Daily'' and ``Weekly'' frequencies apply to site-level inspection and checking activities; ``Monthly'' and ``Quarterly'' to formal documented review of the risk assessment or management arrangement; ``Annual'' to policy, procedure, and legal compliance review. ``Trigger-based'' denotes reviews that must occur in response to a specific event regardless of the calendar cycle (see the Cadence of Assessment sections in each chapter for the specific triggers).

The matrix should be read in conjunction with the full chapter content; the frequency indicated is a minimum, not a maximum. Higher-risk activities, complex sites, and projects where conditions are changing rapidly may require more frequent review.

\begin{landscape}
\small
\begin{longtable}{>{\raggedright\arraybackslash}p{5.8cm}>{\centering\arraybackslash}p{1.0cm}>{\centering\arraybackslash}p{1.0cm}>{\centering\arraybackslash}p{1.2cm}>{\centering\arraybackslash}p{1.2cm}>{\centering\arraybackslash}p{1.0cm}>{\centering\arraybackslash}p{1.8cm}}
\caption{Assessment Cadence Matrix --- \textit{Prudentia Aedificatoria}} \label{tab:cadence} \\
\toprule
\textbf{Assessment Topic} & \textbf{Daily} & \textbf{Weekly} & \textbf{Monthly} & \textbf{Quarterly} & \textbf{Annual} & \textbf{Trigger-based} \\
\midrule
\endfirsthead
\multicolumn{7}{c}{\textit{Table continued}} \\
\toprule
\textbf{Assessment Topic} & \textbf{Daily} & \textbf{Weekly} & \textbf{Monthly} & \textbf{Quarterly} & \textbf{Annual} & \textbf{Trigger-based} \\
\midrule
\endhead
\midrule
\multicolumn{7}{r}{\textit{Continued on next page}} \\
\endfoot
\bottomrule
\endlastfoot
Ch.1 --- Introduction to Construction Risk Assessment & & \checkmark & \checkmark & & \checkmark & \checkmark \\
Ch.2 --- Legal Framework and Duty Holders & & & & \checkmark & \checkmark & \checkmark \\
Ch.3 --- The Risk Assessment Process (Five-Step Method) & & \checkmark & \checkmark & & \checkmark & \checkmark \\
Ch.4 --- Hierarchy of Controls & & \checkmark & \checkmark & & \checkmark & \checkmark \\
Ch.5 --- Manual Handling Risk Assessment & \checkmark & \checkmark & \checkmark & & \checkmark & \checkmark \\
Ch.6 --- Excavation and Groundworks Risk Assessment & \checkmark & \checkmark & \checkmark & & \checkmark & \checkmark \\
Ch.7 --- Working at Height Risk Assessment & \checkmark & \checkmark & \checkmark & & \checkmark & \checkmark \\
Ch.8 --- Scaffolding and Temporary Works Risk Assessment & \checkmark & \checkmark & \checkmark & & \checkmark & \checkmark \\
Ch.9 --- Plant and Mobile Equipment Risk Assessment & \checkmark & \checkmark & \checkmark & & \checkmark & \checkmark \\
Ch.10 --- Noise, Vibration and HAVS Risk Assessment & & \checkmark & \checkmark & \checkmark & \checkmark & \checkmark \\
Ch.11 --- COSHH and Hazardous Substances Risk Assessment & & \checkmark & \checkmark & \checkmark & \checkmark & \checkmark \\
Ch.12 --- Asbestos Management and Risk Assessment & \checkmark & \checkmark & \checkmark & & \checkmark & \checkmark \\
Ch.13 --- Electrical Safety Risk Assessment & \checkmark & \checkmark & \checkmark & & \checkmark & \checkmark \\
Ch.14 --- Fire Risk Assessment & \checkmark & \checkmark & \checkmark & & \checkmark & \checkmark \\
Ch.15 --- Lifting Operations and Equipment Risk Assessment & \checkmark & \checkmark & \checkmark & & \checkmark & \checkmark \\
Ch.16 --- Confined Spaces Risk Assessment & \checkmark & \checkmark & \checkmark & & \checkmark & \checkmark \\
Ch.17 --- Demolition Risk Assessment & \checkmark & \checkmark & \checkmark & & \checkmark & \checkmark \\
Ch.18 --- Temporary Works Risk Assessment & \checkmark & \checkmark & \checkmark & & \checkmark & \checkmark \\
Ch.19 --- Traffic Management and Logistics Risk Assessment & \checkmark & \checkmark & \checkmark & & \checkmark & \checkmark \\
Ch.20 --- Refurbishment and Live Environment Risk Assessment & \checkmark & \checkmark & \checkmark & & \checkmark & \checkmark \\
Ch.21 --- Occupational Health and Dust Risk Assessment & & \checkmark & \checkmark & \checkmark & \checkmark & \checkmark \\
Ch.22 --- Utilities, Services and Ground Conditions Risk Assessment & \checkmark & \checkmark & \checkmark & & \checkmark & \checkmark \\
Ch.23 --- PPE Selection, Maintenance and Provision & & \checkmark & \checkmark & & \checkmark & \checkmark \\
Ch.24 --- Welfare, Site Facilities and Vulnerable Workers Risk Assessment & & \checkmark & \checkmark & & \checkmark & \checkmark \\
Ch.25 --- Mental Health, Stress and Workforce Wellbeing Risk Assessment & & & \checkmark & \checkmark & \checkmark & \checkmark \\
Ch.26 --- Lone Working, Site Security and Violence at Work Risk Assessment & & \checkmark & \checkmark & & \checkmark & \checkmark \\
Ch.27 --- CDM Duties, Documentation and Competence Risk Assessment & & & \checkmark & & \checkmark & \checkmark \\
Ch.28 --- Subcontractor, Supply Chain and Procurement Risk Assessment & & & \checkmark & \checkmark & \checkmark & \checkmark \\
Ch.29 --- Incident Investigation, Reporting and Lessons Learned & & & \checkmark & \checkmark & \checkmark & \checkmark \\
Ch.30 --- Audit, Cadence and Continuous Improvement & & & \checkmark & \checkmark & \checkmark & \checkmark \\
\end{longtable}
\end{landscape}

\section*{Notes on Using the Cadence Matrix}

The Cadence Matrix is a planning reference, not a substitute for the site-specific risk assessment process. The following notes explain how to apply the matrix in practice.

\textbf{Daily frequencies} apply to physical inspection and pre-work checking activities: excavation inspection before operative entry; scaffold inspection before use; daily plant pre-use checks; hot work permit closure; confined space atmospheric monitoring. These daily activities are the front line of risk management and must be embedded in the site routine, not treated as intermittent events.

\textbf{Weekly frequencies} apply to formal documented site safety inspections (typically the site manager's weekly inspection report), welfare facility condition checks, COSHH store inspection, and RIDDOR register update. Weekly inspections must be conducted by a person with sufficient competence to identify the hazards relevant to the current phase of works; the format should prompt the inspector to cover the key risk domains rather than inviting a generic ``all satisfactory'' response.

\textbf{Monthly frequencies} apply to formal review of the risk assessments themselves: checking that the documented controls remain adequate for the current site conditions, that any new activities or hazards introduced during the month are addressed, and that the corrective actions from the previous month's inspection have been implemented. Monthly review is also the appropriate cadence for subcontractor safety performance review, CDM CPP update, and management system leading indicator reporting to the board.

\textbf{Quarterly frequencies} apply to management system-level reviews: health surveillance data review, stress indicator tool analysis, SSIP certificate expiry checks, insurance renewal verification, and the higher-level assessment topics (legal compliance, competence register, and health and safety objectives progress) that do not change month by month but require regular management attention.

\textbf{Annual frequencies} apply to policy and procedure review, legal compliance audit, competence training renewal programme, ISO 45001 certification surveillance audit, and the comprehensive review of each chapter's assessment to ensure it remains current with changes in legislation, HSE guidance, and industry best practice.

\textbf{Trigger-based frequencies} are the most critical category: these are the reviews that must occur in response to a specific event regardless of where that event falls in the calendar cycle. Every chapter identifies the specific triggers that require an immediate reassessment; the most common triggers across all risk domains are: a significant incident or near miss; a change in work method, materials, or team composition; a change in site conditions (weather, groundwater, adjacent structure); a regulatory enforcement action; and a significant programme change.

\section*{Integrating the Cadence Matrix with the Construction Phase Plan}

The Construction Phase Plan (CPP) required by CDM 2015 must describe the arrangements for managing the significant health and safety risks of the project throughout the construction phase. One of the most practical ways to embed the Cadence Matrix into the CPP is to produce a project-specific health and safety monitoring schedule derived from the matrix, identifying which assessment topics are relevant to the project, which frequencies apply, who is responsible for each review, and how the review records will be maintained.

A well-structured CPP health and safety monitoring schedule might, for a complex civil engineering project, include the following site-level activities on a daily basis: excavation pre-entry inspection; pre-lifting check (where cranes and hoists are in use); confined space atmospheric test (where confined space entry is planned); and daily plant pre-use check for all mobile plant. Weekly activities on the same project might include: scaffold inspection by a competent scaffold inspector; site safety inspection by the site manager or safety adviser; welfare facility inspection and cleaning verification; COSHH store condition check; and toolbox talk delivery. Monthly activities might include: full risk assessment review for the current phase of works; subcontractor safety performance review; CDM CPP update; health surveillance review; and corrective action register review. Quarterly activities might include: SSIP certificate expiry audit across the subcontractor list; stress indicator tool data review; insurance certificate currency check; and management system leading indicator dashboard review for the board.

The matrix also provides a framework for the post-project health and safety file assembly. The Health and Safety File required by CDM 2015 must contain enough information to allow future maintenance and construction work to be carried out safely; the assessment topics in the matrix that relate to embedded or residual hazards in the completed building (asbestos, hazardous substances, structural loading, confined spaces, roof access, electrical systems) should inform the content requirements of the HSF and ensure that the relevant risk information is preserved for the benefit of future workers.

\section*{Applying the Cadence Matrix to Different Project Types}

The Cadence Matrix is designed to be read selectively rather than applied in its entirety to every project. The relevant rows will differ significantly between project types:

For a \textbf{major infrastructure project} (motorway, railway, pipeline): the daily and weekly rows for Ch.6 (Excavation), Ch.9 (Plant), Ch.7 (Working at Height), Ch.13 (Electrical Safety), Ch.15 (Lifting), Ch.17 (Demolition where structures are removed), Ch.19 (Traffic Management), and Ch.22 (Utilities) will be of primary importance. Monthly and quarterly reviews for Ch.25 (Mental Health --- accommodation-based long-duration projects present significant stress and fatigue risk), Ch.28 (Supply Chain --- complex multi-tier supply chains on large infrastructure), and Ch.27 (CDM --- complex duty holder arrangements on major projects) will require particular attention.

For a \textbf{high-rise residential or commercial building project}: the daily and weekly rows for Ch.7 (Working at Height), Ch.8 (Scaffolding), Ch.14 (Fire), Ch.13 (Electrical), Ch.15 (Lifting), and Ch.16 (Confined Spaces --- basement and plantroom work) will dominate. CDM (Ch.27) and Supply Chain (Ch.28) require close management on complex, multi-phase building projects.

For a \textbf{refurbishment and fit-out project}: Ch.12 (Asbestos) and Ch.11 (COSHH) are of elevated importance in the daily and weekly regime. Ch.20 (Live Environment), Ch.14 (Fire), and Ch.13 (Electrical) require daily attention in occupied premises. Welfare (Ch.24) requires attention where the project is in a building that the client continues to occupy.

For a \textbf{small domestic extension or alteration}: the CDM domestic client provision means that the contractor takes on the client's CDM duties. The relevant assessment topics are fewer, but the minimum requirements --- welfare (Ch.24), working at height (Ch.7), manual handling (Ch.5), and electrical safety (Ch.13) at a minimum --- still apply with full legal force.
